\documentclass[11pt]{article}

\usepackage[margin=1in]{geometry}
\usepackage{amsmath,amssymb,amsthm}
\usepackage{graphicx}
\usepackage[hidelinks]{hyperref}

\newtheorem{theorem}{Theorem}
\newtheorem{lemma}{Lemma}

\DeclareMathOperator{\SO}{SO}
\DeclareMathOperator{\SU}{SU}
\DeclareMathOperator{\U}{U}
\DeclareMathOperator{\Sp}{Sp}
\DeclareMathOperator{\so}{\mathfrak{so}}
\DeclareMathOperator{\su}{\mathfrak{su}}
\DeclareMathOperator{\rank}{rank}

\title{The Six-Dimensional Case of the Skew-Hermitian Lie Algebra Problem}
\author{Automatic Math Discovery Smoke Test}
\date{June 13, 2026}

\begin{document}
\maketitle

\section*{Source problem}

The source paper asks for the possible dimensions of the Lie algebra
$\mathcal Z(X)$ of skew-hermitian operators when $X$ is a finite-dimensional
real Banach space. It notes that dimensions $3$ and $4$ were known and that
the first open case was dimension $5$.

\begin{center}
\includegraphics[width=\textwidth]{figures/open_problem_crop.png}
\end{center}

\section*{Result}

\begin{theorem}
Let $X$ be a six-dimensional real Banach space. Then
\[
\dim \mathcal Z(X)\in\{0,1,2,3,4,6,7,10,15\}.
\]
Conversely, every value in this set is attained by some six-dimensional real
Banach space.
\end{theorem}

\begin{proof}
By the Auerbach reduction quoted in the source paper, $\mathcal Z(X)$ embeds
as a compact Lie subalgebra of $\so(6)$. Rosenthal's gap theorem says that,
for $\dim X=n$, if
\[
\dim \mathcal Z(X)>\frac{(n-1)(n-2)}{2},
\]
then $X$ is Hilbert and $\dim\mathcal Z(X)=n(n-1)/2$; and equality at
$(n-1)(n-2)/2$ occurs exactly for a non-Euclidean absolute sum of
$\mathbb R$ and a Hilbert space of dimension $n-1$. For $n=6$, this leaves
only the values
\[
0,1,2,3,4,5,6,7,8,9,10,15.
\]
The value $10$ is the equality case in Rosenthal's theorem. It remains to
exclude $5$, $8$, and $9$.

We use two elementary compact Lie algebra facts. First, a compact Lie algebra
is the direct sum of its center and its compact semisimple part. Second, the
only compact simple Lie algebras of dimension at most $9$ are $\su(2)$,
dimension $3$, and $\su(3)$, dimension $8$.

\smallskip
\noindent\textbf{Excluding dimension $5$.}
A five-dimensional compact Lie algebra in $\so(6)$ cannot be abelian, since
an abelian subalgebra of $\so(6)$ has dimension at most $\rank \so(6)=3$.
Thus it must be isomorphic to $\su(2)\oplus\mathbb R^2$.

Consider the associated faithful orthogonal representation on $\mathbb R^6$.
The nontrivial real irreducible $\su(2)$-modules of dimension at most $6$
have dimensions $3$, $4$, and $5$. The central $\mathbb R^2$ must act through
the skew-symmetric commutant of the $\su(2)$ action. A case check gives only
one way for this commutant to have rank at least $2$: the $\su(2)$ module is
a four-dimensional quaternionic irreducible plus two trivial dimensions.
In that case the relevant commutant is $\mathfrak{sp}(1)\oplus\so(2)$, and
the faithful $\mathbb R^2$ uses a circle in the $\mathfrak{sp}(1)$ factor
and the $\so(2)$ on the trivial two-plane.

Hence the connected group generated on the four-dimensional summand is
$\SU(2)$ together with one commuting circle, i.e. $\U(2)$ on $\mathbb C^2$.
This group is transitive on the Euclidean sphere of that four-dimensional
summand. If a norm on $\mathbb R^4\oplus\mathbb R^2$ is invariant under this
action and under the rotation of the two-plane, then for each fixed two-plane
component it depends on the first component only through its Euclidean length,
and for each fixed first component it depends on the two-plane component only
through its Euclidean length. Therefore the norm is invariant under
$\SO(4)\times\SO(2)$, whose Lie algebra has dimension $6+1=7$. Thus a
six-dimensional Banach space cannot have $\dim\mathcal Z(X)=5$.

\smallskip
\noindent\textbf{Excluding dimension $8$.}
Let $\mathfrak g\subset\so(6)$ be compact of dimension $8$. If the center of
$\mathfrak g$ had positive dimension large enough to make the semisimple part
$0$ or $\su(2)$, it would contain an abelian subalgebra of dimension at least
$5$, impossible in $\so(6)$. If the semisimple part were
$\su(2)\oplus\su(2)$, the center would have dimension $2$; but the faithful
six-dimensional orthogonal representations of $\su(2)\oplus\su(2)$ are, up
to adding trivial summands, either the $3+3$ adjoint representation or the
four-dimensional real spin representation of $\so(4)$ plus a two-dimensional
trivial summand. Their skew-symmetric commutants have rank $0$ and $1$,
respectively, so a two-dimensional center cannot act faithfully.

The remaining possibility is $\mathfrak g\cong\su(3)$ acting on $\mathbb C^3$
viewed as $\mathbb R^6$. The group $\SU(3)$ is transitive on the unit sphere
of $\mathbb C^3$. Therefore every $\SU(3)$-invariant norm on $\mathbb R^6$
is Euclidean, and its skew-hermitian Lie algebra is all of $\so(6)$, of
dimension $15$. Thus dimension $8$ is impossible.

\smallskip
\noindent\textbf{Excluding dimension $9$.}
Let $\mathfrak g\subset\so(6)$ be compact of dimension $9$. The possible
semisimple parts are $0$, $\su(2)$, $\su(2)\oplus\su(2)$, $\su(3)$, and
$\su(2)^3$.

The first two require a center of dimension at least $6$, impossible because
abelian subalgebras of $\so(6)$ have dimension at most $3$. The case
$\su(2)\oplus\su(2)$ would require a three-dimensional center, but the
centralizer ranks in all faithful six-dimensional orthogonal representations
listed above are at most $1$.

For the $\su(2)^3$ case, no faithful real orthogonal representation has
dimension $6$: a nontrivial real irreducible of one $\su(2)$ has dimension at
least $3$, the smallest faithful module for two factors has real dimension
$4$, and the smallest way to include the third factor already has dimension
at least $4+3=7$.

Thus the only remaining case is $\su(3)\oplus\mathbb R$, acting as
$\mathfrak u(3)$ on $\mathbb C^3$. The group $\U(3)$ is transitive on the
unit sphere of $\mathbb C^3$, so again any invariant norm is Euclidean and
the skew-hermitian Lie algebra enlarges to $\so(6)$, dimension $15$.
Therefore dimension $9$ is impossible.

\smallskip
We now realize the remaining values. Let $H_k$ denote the real Hilbert space
of dimension $k$. For a finite $\ell_\infty$-sum of Euclidean blocks, the
identity component of the isometry group is the product of the orthogonal
groups on the Euclidean blocks, so the Lie algebra dimension is
$\sum_j k_j(k_j-1)/2$. Thus the following six-dimensional spaces realize the
announced values:
\[
\begin{array}{c|c}
\dim\mathcal Z(X) & X\\ \hline
0 & \ell_\infty^6\\
1 & H_2\oplus_\infty \ell_\infty^4\\
2 & H_2\oplus_\infty H_2\oplus_\infty \ell_\infty^2\\
3 & H_3\oplus_\infty \ell_\infty^3\\
4 & H_3\oplus_\infty H_2\oplus_\infty \mathbb R\\
6 & H_4\oplus_\infty \ell_\infty^2\\
7 & H_4\oplus_\infty H_2\\
10 & H_5\oplus_\infty \mathbb R\\
15 & H_6.
\end{array}
\]
This proves both necessity and realization.
\end{proof}

\section*{Verification notes}

The script \texttt{code/check\_six\_dimension\_values.py} lists the
Euclidean-block construction values and confirms that Rosenthal's theorem
reduces the finite exclusion list to $5$, $8$, and $9$. It is only a sanity
check; the proof of the exclusions is the compact Lie algebra argument above.

\section*{Novelty check}

Before packaging this candidate partial result, I searched the local parsed
corpus and run artifacts for the exact source arXiv id and phrases such as
\texttt{dim Z(X)}, \texttt{dimension of mathcal-Z(X)},
\texttt{six-dimensional real Banach}, and \texttt{skew-hermitian}. The only
nearby results found in this run were the source problem statement and the
five-dimensional packet.

I also ran bounded web searches on 2026-06-13 using exact and near-exact
phrases around ``dim Z(X)'', ``six-dimensional real Banach space'',
``skew-hermitian operators'', and Rosenthal's Lie algebra paper. These searches
found adjacent numerical-index and skew-hermitian-operator literature, but no
direct later classification of the six-dimensional case.

\section*{Scope}

This is a candidate solution of the six-dimensional subcase of the general
possible-values problem. It does not solve the arbitrary-dimensional problem.

\end{document}
